\begin{frame}
	\frametitle{Esparsidade: onde mora a eficiência}
	\only<1>{
		\begin{columns}
			\column{.55\linewidth}
			\begin{center}
				\includegraphics[height=3.3cm]{../../00-thesis/monography/images/sparsity}\par
				{\tiny Eixo horizontal: tempo. Eixo vertical: índice do neurônio. Fonte: elaborado pelo autor.}
			\end{center}
			\column{.42\linewidth}
			\small
			\begin{itemize}
				\item Na maior parte do tempo, a \textbf{grande maioria} dos neurônios está silenciosa.
				\item Pulsos são representados como \textbf{uns} dispersamente distribuídos em uma sequência de \textbf{zeros} \cite{eshraghian2023training}.
				\item Um neurônio LIF bem parametrizado tolera ruído: age como um \textbf{filtro passa-baixa}, só disparando quando o sinal se acumula de forma sustentada.
			\end{itemize}
		\end{columns}
		% [SOFTWARE] SNN -> Codificacao Poisson (imagem) | t ~ 15/29,
		% aponte que cada pixel e um "neuronio" sorteando disparo
		\abrirNoSoftware{SNN -> Codificação Poisson (imagem)}{snn.poisson_image}
	}
	\only<2>{
		\begin{alertblock}{A promessa, com precisão}
			Esparsidade só vira economia de energia real em hardware que sabe \textbf{pular} o zero. Em GPU/TPU convencional, a esparsidade é invisível ao hardware --- ele computa a multiplicação por zero do mesmo jeito.
		\end{alertblock}
		\vspace{4pt}
		\par Volta-se a este ponto na próxima seção, sobre hardware neuromórfico.
	}
\end{frame}
